\documentclass[12pt,a4paper]{article}

% --- Nevyhnutné balíky ---
\usepackage[utf8]{inputenc}
\usepackage[slovak]{babel}
\usepackage{float}       % Rieši chybu s [H]
\usepackage{pgfplots}    % Rieši grafy
\usepackage{geometry}    % Nastavenie okrajov
\geometry{margin=2cm}
\pgfplotsset{compat=1.18}

\begin{document}

\section{Detailné výsledky testovacích scenárov}

\subsection{Test 1: Základné nastavenie systému}
% ═══════════════════════════════════════════════════════════════════════════════

V prvom testovacom scenári overujeme správanie algoritmu pri vyváženom nastavení váh. Tento test slúži ako referenčný bod (baseline) pre ďalšie simulácie. Váhy sú rozdelené tak, aby hlavná kategória mala miernu prioritu (20 \%), zatiaľ čo ostatné faktory sú zastúpené rovnomerne (16 \%), ako uvádza tabuľka \ref{tab:vahy_test1}.

\begin{table}[H]
\centering
\begin{tabular}{|l|c|}
\hline
\textbf{Faktor} & \textbf{Váha ($w_i$)} \\
\hline
Hlavná kategória ($w_1$) & 20.00\% \\
Podkategória ($w_2$) & 16.00\% \\
Vzdialenosť ($w_3$) & 16.00\% \\
Rating ($w_4$) & 16.00\% \\
Popularita ($w_5$) & 16.00\% \\
Obľúbený bonus ($w_6$) & 16.00\% \\
\hline
\textbf{Súčet} & \textbf{100\%} \\
\hline
\end{tabular}
\caption{Konfigurácia váh pre Test 1}
\label{tab:vahy_test1}
\end{table}

Pre tento scenár boli vygenerované odporúčania pre používateľa so záujmom o kategóriu \textit{Sports}. V tabuľke \ref{tab:data_test1} sú zobrazené základné atribúty piatich najúspešnejších podujatí pred procesom normalizácie.

\begin{table}[H]
\centering
\small
\begin{tabular}{|c|l|l|l|c|c|c|}
\hline
\textbf{\#} & \textbf{Názov udalosti} & \textbf{Kategória} & \textbf{Podkategória} & \textbf{Vzdial.} & \textbf{Rating} & \textbf{Popularita} \\
\hline
1 & Udalosť 5 & Sports    & Strength \& Cond. & 13.36 km & 4.9 (13) & 9/16 \\
2 & Udalosť 15 & Sports    & Strength \& Cond. & 11.73 km & 4.3 (17) & 4/71 \\
3 & Udalosť 7 & Sports    & Strength \& Cond. & 15.63 km & 4.2 (14) & 6/69 \\
4 & Udalosť 2 & Education & Jazykový kurz     & 1.37 km  & 4.7 (10) & 5/7 \\
5 & Udalosť 4 & Sports    & Team sports       & 15.80 km & 4.9 (1)  & 6/9 \\
\hline
\end{tabular}
\caption{Vstupné údaje pre TOP 5 odporúčaní - Test 1}
\label{tab:data_test1}
\end{table}

V ďalšom kroku algoritmus transformuje tieto údaje na bezrozmerné normalizované hodnoty v rozsahu $\langle 0, 1 \rangle$. Tieto hodnoty (tabuľka \ref{tab:norm_test1}) následne vstupujú do procesu váženého súčtu.

\begin{table}[H]
\centering
\small
\begin{tabular}{|c|l|c|c|c|c|c|c|}
\hline
\textbf{\#} & \textbf{Event} & $C_{main}$ & $C_{sub}$ & $D$ & $R$ & $P$ & $B_{fav}$ \\
\hline
1 & Udalosť 5 & 0.500 & 0.125 & 0.428 & 0.98 & 0.563 & 1.0 \\
2 & Udalosť 15 & 0.500 & 0.125 & 0.460 & 0.86 & 0.056 & 1.0 \\
3 & Udalosť 7 & 0.500 & 0.125 & 0.390 & 0.84 & 0.087 & 1.0 \\
4 & Udalosť 2 & 0.375 & 0.000 & 0.879 & 0.94 & 0.714 & 0.0 \\
5 & Udalosť 4 & 0.500 & 0.125 & 0.387 & 0.98 & 0.667 & 0.0 \\
\hline
\end{tabular}
\caption{Normalizované hodnoty pred aplikáciou váh - Test 1}
\label{tab:norm_test1}
\end{table}

Finálny výsledok a vplyv jednotlivých váh na celkové skóre sú zobrazené v tabuľke \ref{tab:vysledky_test1}. Najvyššie skóre (59,7 \%) dosiahla Udalosť 5, predovšetkým vďaka kombinácii vysokého ratingu a bonusu za obľúbenú kategóriu.

\begin{table}[H]
\centering
\small
\begin{tabular}{|c|l|c|c|c|c|c|c|c|}
\hline
\textbf{\#} & \textbf{Event} & \textbf{Score} & $w_1 C$ & $w_2 C_s$ & $w_3 D$ & $w_4 R$ & $w_5 P$ & $w_6 B$ \\
\hline
1 & Udalosť 5  & 59.7\% & 10.0\% & 2.0\% & 6.8\%  & 15.8\% & 9.0\%  & 16.0\% \\
2 & Udalosť 15 & 50.1\% & 10.0\% & 2.0\% & 7.4\%  & 13.9\% & 0.9\%  & 16.0\% \\
3 & Udalosť 7  & 49.0\% & 10.0\% & 2.0\% & 6.2\%  & 13.3\% & 1.4\%  & 16.0\% \\
4 & Udalosť 2  & 47.9\% & 7.5\%  & 0.0\% & 14.1\% & 14.9\% & 11.4\% & 0.0\% \\
5 & Udalosť 4  & 44.7\% & 10.0\% & 2.0\% & 6.2\%  & 15.8\% & 10.7\% & 0.0\% \\
\hline
\end{tabular}
\caption{Rozbor príspevkov k finálnemu skóre - Test 1}
\label{tab:vysledky_test1}
\end{table}

Nasledujúce vizualizácie graficky znázorňujú štruktúru skóre pre všetkých päť odporúčaných udalostí. To umožňuje rýchlu identifikáciu faktorov, ktoré najviac ovplyvnili výsledné poradie.

\begin{figure}[H]
\centering
\begin{minipage}{0.48\textwidth}
\centering
\begin{tikzpicture}
\begin{axis}[
    xbar, width=6cm, height=5cm,
    xlabel={Príspevok (\%)},
    symbolic y coords={Bonus, Popularita, Rating, Vzdialenosť, Podkat., Hl. kat.},
    ytick=data, nodes near coords, nodes near coords align={horizontal},
    title={Udalosť 5 - 59.7\%},
    xmin=0, xmax=25, bar width=0.25cm,
]
\addplot coordinates {(10.0,Hl. kat.) (2.0,Podkat.) (6.8,Vzdialenosť) (15.8,Rating) (9.0,Popularita) (16.0,Bonus)};
\end{axis}
\end{tikzpicture}
\end{minipage}
\hfill
\begin{minipage}{0.48\textwidth}
\centering
\begin{tikzpicture}
\begin{axis}[
    xbar, width=6cm, height=5cm,
    xlabel={Príspevok (\%)},
    symbolic y coords={Bonus, Popularita, Rating, Vzdialenosť, Podkat., Hl. kat.},
    ytick=data, nodes near coords, nodes near coords align={horizontal},
    title={Udalosť 15 - 50.1\%},
    xmin=0, xmax=25, bar width=0.25cm,
]
\addplot coordinates {(10.0,Hl. kat.) (2.0,Podkat.) (7.4,Vzdialenosť) (13.9,Rating) (0.9,Popularita) (16.0,Bonus)};
\end{axis}
\end{tikzpicture}
\end{minipage}
\caption{Vizuálny rozbor skóre pre Test 1 (Udalosť 5 a 15)}
\end{figure}

\begin{figure}[H]
\centering
\begin{minipage}{0.48\textwidth}
\centering
\begin{tikzpicture}
\begin{axis}[
    xbar, width=6cm, height=5cm,
    xlabel={Príspevok (\%)},
    symbolic y coords={Bonus, Popularita, Rating, Vzdialenosť, Podkat., Hl. kat.},
    ytick=data, nodes near coords, nodes near coords align={horizontal},
    title={Udalosť 7 - 49.0\%},
    xmin=0, xmax=25, bar width=0.25cm,
]
\addplot coordinates {(10.0,Hl. kat.) (2.0,Podkat.) (6.2,Vzdialenosť) (13.3,Rating) (1.4,Popularita) (16.0,Bonus)};
\end{axis}
\end{tikzpicture}
\end{minipage}
\hfill
\begin{minipage}{0.48\textwidth}
\centering
\begin{tikzpicture}
\begin{axis}[
    xbar, width=6cm, height=5cm,
    xlabel={Príspevok (\%)},
    symbolic y coords={Bonus, Popularita, Rating, Vzdialenosť, Podkat., Hl. kat.},
    ytick=data, nodes near coords, nodes near coords align={horizontal},
    title={Udalosť 2 - 47.9\%},
    xmin=0, xmax=25, bar width=0.25cm,
]
\addplot coordinates {(7.5,Hl. kat.) (0.0,Podkat.) (14.1,Vzdialenosť) (14.9,Rating) (11.4,Popularita) (0.0,Bonus)};
\end{axis}
\end{tikzpicture}
\end{minipage}
\caption{Vizuálny rozbor skóre pre Test 1 (Udalosť 7 a 2)}
\end{figure}

\begin{figure}[H]
\centering
\begin{tikzpicture}
\begin{axis}[
    xbar, width=10cm, height=5cm,
    xlabel={Príspevok k skóre (\%)},
    symbolic y coords={Bonus, Popularita, Rating, Vzdialenosť, Podkat., Hl. kat.},
    ytick=data, nodes near coords, nodes near coords align={horizontal},
    title={Udalosť 4 - 44.7\%},
    xmin=0, xmax=25, bar width=0.4cm,
]
\addplot coordinates {(10.0,Hl. kat.) (2.0,Podkat.) (6.2,Vzdialenosť) (15.8,Rating) (10.7,Popularita) (0.0,Bonus)};
\end{axis}
\end{tikzpicture}
\caption{Vizuálny rozbor skóre pre Test 1 (Udalosť 4)}
\end{figure}


% ═══════════════════════════════════════════════════════════════════════════════
\subsection{Test 2: Prioritizácia hlavnej kategórie (Hl. kategória 60\%)}
% ═══════════════════════════════════════════════════════════════════════════════

V druhom scenári simulujeme správanie systému pre používateľa, ktorý má vyhranený záujem o konkrétnu oblasť (v tomto prípade \textit{Sports}) a ostatné faktory ako vzdialenosť či popularita sú preňho druhoradé. Váha hlavnej kategórie $w_1$ bola zvýšená na 60 \%, zatiaľ čo vplyv ostatných zložiek bol redukovaný na 8 \%, ako je uvedené v tabuľke \ref{tab:vahy_test2}.

\begin{table}[H]
\centering
\begin{tabular}{|l|c|}
\hline
\textbf{Faktor} & \textbf{Váha ($w_i$)} \\
\hline
Hlavná kategória ($w_1$) & 60.00\% \\
Podkategória ($w_2$) & 8.00\% \\
Vzdialenosť ($w_3$) & 8.00\% \\
Rating ($w_4$) & 8.00\% \\
Popularita ($w_5$) & 8.00\% \\
Obľúbený bonus ($w_6$) & 8.00\% \\
\hline
\textbf{Súčet} & \textbf{100\%} \\
\hline
\end{tabular}
\caption{Konfigurácia váh pre Test 2}
\label{tab:vahy_test2}
\end{table}

Vstupné údaje pre tento test (tabuľka \ref{tab:data_test2}) ukazujú, že všetkých TOP 5 výsledkov patrí do kategórie \textit{Sports}. Systém v tomto prípade filtruje podujatia primárne podľa tematickej zhody.

\begin{table}[H]
\centering
\small
\begin{tabular}{|c|l|l|l|c|c|c|}
\hline
\textbf{\#} & \textbf{Názov udalosti} & \textbf{Kategória} & \textbf{Podkategória} & \textbf{Vzdial.} & \textbf{Rating} & \textbf{Popularita} \\
\hline
1 & Udalosť 5  & Sports & Strength \& Cond. & 13.36 km & 4.9 (13) & 9/16 \\
2 & Udalosť 15 & Sports & Strength \& Cond. & 11.73 km & 4.3 (17) & 4/71 \\
3 & Udalosť 7  & Sports & Strength \& Cond. & 15.63 km & 4.2 (14) & 6/69 \\
4 & Udalosť 4  & Sports & Team sports       & 15.80 km & 4.9 (1)  & 6/9 \\
5 & Udalosť 1  & Sports & Team sports       & 8.76 km  & 4.6 (16) & 5/9 \\
\hline
\end{tabular}
\caption{Vstupné údaje pre TOP 5 odporúčaní - Test 2}
\label{tab:data_test2}
\end{table}

Po normalizácii údajov (tabuľka \ref{tab:norm_test2}) vidíme, že zhoda v hlavnej kategórii ($C_{main}$) dosahuje u všetkých subjektov rovnakú hodnotu, čo pri vysokej váhe $w_1$ vedie k veľmi vyrovnaným výsledkom.

\begin{table}[H]
\centering
\small
\begin{tabular}{|c|l|c|c|c|c|c|c|}
\hline
\textbf{\#} & \textbf{Event} & $C_{main}$ & $C_{sub}$ & $D$ & $R$ & $P$ & $B_{fav}$ \\
\hline
1 & Udalosť 5  & 0.500 & 0.125 & 0.428 & 0.98 & 0.563 & 1.0 \\
2 & Udalosť 15 & 0.500 & 0.125 & 0.460 & 0.86 & 0.056 & 1.0 \\
3 & Udalosť 7  & 0.500 & 0.125 & 0.390 & 0.84 & 0.087 & 1.0 \\
4 & Udalosť 4  & 0.500 & 0.125 & 0.387 & 0.98 & 0.667 & 0.0 \\
5 & Udalosť 1  & 0.500 & 0.125 & 0.533 & 0.92 & 0.556 & 0.0 \\
\hline
\end{tabular}
\caption{Normalizované hodnoty - Test 2}
\label{tab:norm_test2}
\end{table}

Finálne skóre v tabuľke \ref{tab:vysledky_test2} potvrdzuje dominanciu kategórie. Udalosti, ktoré majú navyše priradený bonus za obľúbenú položku ($w_6 B$), získali signifikantný náskok pred ostatnými športovými podujatiami.

\begin{table}[H]
\centering
\small
\begin{tabular}{|c|l|c|c|c|c|c|c|c|}
\hline
\textbf{\#} & \textbf{Event} & \textbf{Score} & $w_1 C$ & $w_2 C_s$ & $w_3 D$ & $w_4 R$ & $w_5 P$ & $w_6 B$ \\
\hline
1 & Udalosť 5  & 54.8\% & 30.0\% & 1.0\% & 3.4\% & 7.9\% & 4.5\% & 8.0\% \\
2 & Udalosť 15 & 50.1\% & 30.0\% & 1.0\% & 3.7\% & 6.9\% & 0.5\% & 8.0\% \\
3 & Udalosť 7  & 49.5\% & 30.0\% & 1.0\% & 3.1\% & 6.7\% & 0.7\% & 8.0\% \\
4 & Udalosť 4  & 47.3\% & 30.0\% & 1.0\% & 3.1\% & 7.9\% & 5.3\% & 0.0\% \\
5 & Udalosť 1  & 47.1\% & 30.0\% & 1.0\% & 4.3\% & 7.4\% & 4.4\% & 0.0\% \\
\hline
\end{tabular}
\caption{Rozbor príspevkov k finálnemu skóre - Test 2}
\label{tab:vysledky_test2}
\end{table}

Vizualizácia rozdelenia skóre na obrázku nižšie jasne demonštruje, že hlavná kategória tvorí u všetkých odporúčaní najväčšiu časť výslednej hodnoty.

\begin{figure}[H]
\centering
\begin{minipage}{0.48\textwidth}
\centering
\begin{tikzpicture}
\begin{axis}[
    xbar, width=6cm, height=5cm,
    xlabel={Príspevok (\%)},
    symbolic y coords={Bonus, Popularita, Rating, Vzdialenosť, Podkat., Hl. kat.},
    ytick=data, nodes near coords, nodes near coords align={horizontal},
    title={Udalosť 5 - 54.8\%},
    xmin=0, xmax=35, bar width=0.25cm,
]
\addplot coordinates {(30.0,Hl. kat.) (1.0,Podkat.) (3.4,Vzdialenosť) (7.9,Rating) (4.5,Popularita) (8.0,Bonus)};
\end{axis}
\end{tikzpicture}
\end{minipage}
\hfill
\begin{minipage}{0.48\textwidth}
\centering
\begin{tikzpicture}
\begin{axis}[
    xbar, width=6cm, height=5cm,
    xlabel={Príspevok (\%)},
    symbolic y coords={Bonus, Popularita, Rating, Vzdialenosť, Podkat., Hl. kat.},
    ytick=data, nodes near coords, nodes near coords align={horizontal},
    title={Udalosť 15 - 50.1\%},
    xmin=0, xmax=35, bar width=0.25cm,
]
\addplot coordinates {(30.0,Hl. kat.) (1.0,Podkat.) (3.7,Vzdialenosť) (6.9,Rating) (0.5,Popularita) (8.0,Bonus)};
\end{axis}
\end{tikzpicture}
\end{minipage}
\caption{Vizuálny rozbor skóre pre Test 2 (Udalosť 5 a 15)}
\end{figure}

\begin{figure}[H]
\centering
\begin{minipage}{0.48\textwidth}
\centering
\begin{tikzpicture}
\begin{axis}[
    xbar, width=6cm, height=5cm,
    xlabel={Príspevok (\%)},
    symbolic y coords={Bonus, Popularita, Rating, Vzdialenosť, Podkat., Hl. kat.},
    ytick=data, nodes near coords, nodes near coords align={horizontal},
    title={Udalosť 7 - 49.5\%},
    xmin=0, xmax=35, bar width=0.25cm,
]
\addplot coordinates {(30.0,Hl. kat.) (1.0,Podkat.) (3.1,Vzdialenosť) (6.7,Rating) (0.7,Popularita) (8.0,Bonus)};
\end{axis}
\end{tikzpicture}
\end{minipage}
\hfill
\begin{minipage}{0.48\textwidth}
\centering
\begin{tikzpicture}
\begin{axis}[
    xbar, width=6cm, height=5cm,
    xlabel={Príspevok (\%)},
    symbolic y coords={Bonus, Popularita, Rating, Vzdialenosť, Podkat., Hl. kat.},
    ytick=data, nodes near coords, nodes near coords align={horizontal},
    title={Udalosť 4 - 47.3\%},
    xmin=0, xmax=35, bar width=0.25cm,
]
\addplot coordinates {(30.0,Hl. kat.) (1.0,Podkat.) (3.1,Vzdialenosť) (7.9,Rating) (5.3,Popularita) (0.0,Bonus)};
\end{axis}
\end{tikzpicture}
\end{minipage}
\caption{Vizuálny rozbor skóre pre Test 2 (Udalosť 7 a 4)}
\end{figure}

\begin{figure}[H]
\centering
\begin{tikzpicture}
\begin{axis}[
    xbar, width=10cm, height=5cm,
    xlabel={Príspevok k skóre (\%)},
    symbolic y coords={Bonus, Popularita, Rating, Vzdialenosť, Podkat., Hl. kat.},
    ytick=data, nodes near coords, nodes near coords align={horizontal},
    title={Udalosť 1 - 47.1\%},
    xmin=0, xmax=35, bar width=0.4cm,
]
\addplot coordinates {(30.0,Hl. kat.) (1.0,Podkat.) (4.3,Vzdialenosť) (7.4,Rating) (4.4,Popularita) (0.0,Bonus)};
\end{axis}
\end{tikzpicture}
\caption{Vizuálny rozbor skóre pre Test 2 (Udalosť 1)}
\end{figure}


% ═══════════════════════════════════════════════════════════════════════════════
\subsection{Test 3: Prioritizácia špecifickej podkategórie (Podkategória 60 \%)}
% ═══════════════════════════════════════════════════════════════════════════════

Tretí scenár sa zameriava na situáciu, kedy má používateľ veľmi konkrétnu preferenciu v rámci danej oblasti, napríklad hľadá špecifický typ tréningu alebo vzdelávania. Váha faktora $w_2$ (Podkategória) bola navýšená na 60 \%, zatiaľ čo ostatné atribúty majú rovnomerný vplyv 8 \%, ako je znázornené v tabuľke \ref{tab:vahy_test3}.

\begin{table}[H]
\centering
\begin{tabular}{|l|c|}
\hline
\textbf{Faktor} & \textbf{Váha ($w_i$)} \\
\hline
Hlavná kategória ($w_1$) & 8.00\% \\
Podkategória ($w_2$) & 60.00\% \\
Vzdialenosť ($w_3$) & 8.00\% \\
Rating ($w_4$) & 8.00\% \\
Popularita ($w_5$) & 8.00\% \\
Obľúbený bonus ($w_6$) & 8.00\% \\
\hline
\textbf{Súčet} & \textbf{100\%} \\
\hline
\end{tabular}
\caption{Nastavenie váh algoritmu pre Test 3}
\label{tab:vahy_test3}
\end{table}

Vstupné údaje v tabuľke \ref{tab:data_test3} ukazujú, že do výberu TOP 5 sa dostali udalosti z kategórií \textit{Sports} a \textit{Education}, ktoré najlepšie zodpovedajú hľadanej podkategórii.

\begin{table}[H]
\centering
\small
\begin{tabular}{|c|l|l|l|c|c|c|}
\hline
\textbf{\#} & \textbf{Názov udalosti} & \textbf{Kategória} & \textbf{Podkategória} & \textbf{Vzdial.} & \textbf{Rating} & \textbf{Popularita} \\
\hline
1 & Udalosť 5  & Sports    & Strength \& Cond. & 13.36 km & 4.9 (13) & 9/16 \\
2 & Udalosť 15 & Sports    & Strength \& Cond. & 11.73 km & 4.3 (17) & 4/71 \\
3 & Udalosť 7  & Sports    & Strength \& Cond. & 15.63 km & 4.2 (14) & 6/69 \\
4 & Udalosť 2  & Education & Osobný rozvoj     & 21.27 km & 4.2 (17) & 15/63 \\
5 & Udalosť 4  & Education & Osobný rozvoj     & 12.08 km & 4.3 (15) & 6/83 \\
\hline
\end{tabular}
\caption{Vstupné údaje pre TOP 5 udalostí - Test 3}
\label{tab:data_test3}
\end{table}



Po procese normalizácie (tabuľka \ref{tab:norm_test3}) systém vypočítal príspevky jednotlivých faktorov k celkovému skóre. Všimnite si nárast vplyvu $w_2 C_s$ pri udalostiach 4 a 5.

\begin{table}[H]
\centering
\small
\begin{tabular}{|c|l|c|c|c|c|c|c|}
\hline
\textbf{\#} & \textbf{Event} & $C_{main}$ & $C_{sub}$ & $D$ & $R$ & $P$ & $B_{fav}$ \\
\hline
1 & Udalosť 5  & 0.500 & 0.125 & 0.428 & 0.98 & 0.563 & 1.0 \\
2 & Udalosť 15 & 0.500 & 0.125 & 0.460 & 0.86 & 0.056 & 1.0 \\
3 & Udalosť 7  & 0.500 & 0.125 & 0.390 & 0.84 & 0.087 & 1.0 \\
4 & Udalosť 2  & 0.375 & 0.250 & 0.320 & 0.84 & 0.238 & 0.0 \\
5 & Udalosť 4  & 0.375 & 0.250 & 0.453 & 0.86 & 0.072 & 0.0 \\
\hline
\end{tabular}
\caption{Normalizované hodnoty pred aplikáciou váh - Test 3}
\label{tab:norm_test3}
\end{table}

Výsledné hodnotenie (tabuľka \ref{tab:vysledky_test3}) dokumentuje, že pri takto nastavenom systéme dokáže silná zhoda v podkategórii konkurovať aj podujatiam s vyšším celkovým ratingom či bonusom za obľúbenosť.

\begin{table}[H]
\centering
\small
\begin{tabular}{|c|l|c|c|c|c|c|c|c|}
\hline
\textbf{\#} & \textbf{Event} & \textbf{Score} & $w_1 C$ & $w_2 C_s$ & $w_3 D$ & $w_4 R$ & $w_5 P$ & $w_6 B$ \\
\hline
1 & Udalosť 5  & 35.3\% & 4.0\% & 7.5\%  & 3.4\% & 7.9\% & 4.5\% & 8.0\% \\
2 & Udalosť 15 & 30.6\% & 4.0\% & 7.5\%  & 3.7\% & 6.9\% & 0.5\% & 8.0\% \\
3 & Udalosť 7  & 30.0\% & 4.0\% & 7.5\%  & 3.1\% & 6.7\% & 0.7\% & 8.0\% \\
4 & Udalosť 2  & 29.2\% & 3.0\% & 15.0\% & 2.6\% & 6.7\% & 1.9\% & 0.0\% \\
5 & Udalosť 4  & 29.2\% & 3.0\% & 15.0\% & 3.6\% & 7.0\% & 0.6\% & 0.0\% \\
\hline
\end{tabular}
\caption{Detailný rozbor príspevkov k finálnemu skóre - Test 3}
\label{tab:vysledky_test3}
\end{table}

Nasledujúce vizualizácie graficky potvrdzujú, že podkategória (druhý stĺpec zospodu) tvorí pri udalostiach 4 a 5 najvýraznejšiu zložku ich odporúčania.

\begin{figure}[H]
\centering
\begin{minipage}{0.48\textwidth}
\centering
\begin{tikzpicture}
\begin{axis}[
    xbar, width=6cm, height=5cm,
    xlabel={Príspevok (\%)},
    symbolic y coords={Bonus, Popularita, Rating, Vzdialenosť, Podkat., Hl. kat.},
    ytick=data, nodes near coords, nodes near coords align={horizontal},
    title={Udalosť 5 - 35.3\%},
    xmin=0, xmax=20, bar width=0.25cm,
]
\addplot coordinates {(4.0,Hl. kat.) (7.5,Podkat.) (3.4,Vzdialenosť) (7.9,Rating) (4.5,Popularita) (8.0,Bonus)};
\end{axis}
\end{tikzpicture}
\end{minipage}
\hfill
\begin{minipage}{0.48\textwidth}
\centering
\begin{tikzpicture}
\begin{axis}[
    xbar, width=6cm, height=5cm,
    xlabel={Príspevok (\%)},
    symbolic y coords={Bonus, Popularita, Rating, Vzdialenosť, Podkat., Hl. kat.},
    ytick=data, nodes near coords, nodes near coords align={horizontal},
    title={Udalosť 15 - 30.6\%},
    xmin=0, xmax=20, bar width=0.25cm,
]
\addplot coordinates {(4.0,Hl. kat.) (7.5,Podkat.) (3.7,Vzdialenosť) (6.9,Rating) (0.5,Popularita) (8.0,Bonus)};
\end{axis}
\end{tikzpicture}
\end{minipage}
\caption{Breakdown skóre – Test 3, Udalosť 5 a 15}
\end{figure}

\begin{figure}[H]
\centering
\begin{minipage}{0.48\textwidth}
\centering
\begin{tikzpicture}
\begin{axis}[
    xbar, width=6cm, height=5cm,
    xlabel={Príspevok (\%)},
    symbolic y coords={Bonus, Popularita, Rating, Vzdialenosť, Podkat., Hl. kat.},
    ytick=data, nodes near coords, nodes near coords align={horizontal},
    title={Udalosť 7 - 30.0\%},
    xmin=0, xmax=20, bar width=0.25cm,
]
\addplot coordinates {(4.0,Hl. kat.) (7.5,Podkat.) (3.1,Vzdialenosť) (6.7,Rating) (0.7,Popularita) (8.0,Bonus)};
\end{axis}
\end{tikzpicture}
\end{minipage}
\hfill
\begin{minipage}{0.48\textwidth}
\centering
\begin{tikzpicture}
\begin{axis}[
    xbar, width=6cm, height=5cm,
    xlabel={Príspevok (\%)},
    symbolic y coords={Bonus, Popularita, Rating, Vzdialenosť, Podkat., Hl. kat.},
    ytick=data, nodes near coords, nodes near coords align={horizontal},
    title={Udalosť 2 - 29.2\%},
    xmin=0, xmax=20, bar width=0.25cm,
]
\addplot coordinates {(3.0,Hl. kat.) (15.0,Podkat.) (2.6,Vzdialenosť) (6.7,Rating) (1.9,Popularita) (0.0,Bonus)};
\end{axis}
\end{tikzpicture}
\end{minipage}
\caption{Breakdown skóre – Test 3, Udalosť 7 a 2}
\end{figure}

\begin{figure}[H]
\centering
\begin{tikzpicture}
\begin{axis}[
    xbar, width=10cm, height=5cm,
    xlabel={Príspevok k skóre (\%)},
    symbolic y coords={Bonus, Popularita, Rating, Vzdialenosť, Podkat., Hl. kat.},
    ytick=data, nodes near coords, nodes near coords align={horizontal},
    title={Udalosť 4 - 29.2\%},
    xmin=0, xmax=20, bar width=0.4cm,
]
\addplot coordinates {(3.0,Hl. kat.) (15.0,Podkat.) (3.6,Vzdialenosť) (7.0,Rating) (0.6,Popularita) (0.0,Bonus)};
\end{axis}
\end{tikzpicture}
\caption{Breakdown skóre - Test 3, Udalosť 4}
\end{figure}


% ═══════════════════════════════════════════════════════════════════════════════
\subsection{Test 4: Prioritizácia geografickej dostupnosti (Vzdialenosť 60 \%)}
% ═══════════════════════════════════════════════════════════════════════════════

Štvrtý testovací scenár simuluje správanie systému pre používateľa, ktorého primárnym kritériom je čo najlepšia dostupnosť podujatia. Váha faktora $w_3$ (Vzdialenosť) bola zvýšená na 60 \%, čo výrazne znevýhodňuje vzdialenejšie podujatia bez ohľadu na ich kvalitu či tematickú zhodu. Konfigurácia je v tabuľke \ref{tab:vahy_test4}.

\begin{table}[H]
\centering
\begin{tabular}{|l|c|}
\hline
\textbf{Faktor} & \textbf{Váha ($w_i$)} \\
\hline
Hlavná kategória ($w_1$) & 8.00\% \\
Podkategória ($w_2$) & 8.00\% \\
Vzdialenosť ($w_3$) & 60.00\% \\
Rating ($w_4$) & 8.00\% \\
Popularita ($w_5$) & 8.00\% \\
Obľúbený bonus ($w_6$) & 8.00\% \\
\hline
\textbf{Súčet} & \textbf{100\%} \\
\hline
\end{tabular}
\caption{Konfigurácia váh pre Test 4}
\label{tab:vahy_test4}
\end{table}

Vstupný súbor dát (tabuľka \ref{tab:data_test4}) jasne ukazuje, že všetky odporúčané udalosti sa nachádzajú v tesnej blízkosti používateľa (do 3 km). Do výberu sa tak dostala aj kategória \textit{Family \& Kids}, ktorá v predchádzajúcich testoch absentovala.

\begin{table}[H]
\centering
\small
\begin{tabular}{|c|l|l|l|c|c|c|}
\hline
\textbf{\#} & \textbf{Názov udalosti} & \textbf{Kategória} & \textbf{Podkategória} & \textbf{Vzdial.} & \textbf{Rating} & \textbf{Popularita} \\
\hline
1 & Udalosť 2A & Education      & Jazykový kurz   & 1.37 km & 4.7 (10) & 5/7 \\
2 & Udalosť 5  & Education      & Vzdelávací kurz & 1.59 km & 4.4 (11) & 5/34 \\
3 & Udalosť 1  & Family \& Kids & Rodinná akcia   & 1.05 km & 4.7 (2)  & 6/99 \\
4 & Udalosť 8  & Sports         & Aquatic sports  & 2.54 km & 3.4 (4)  & 10/34 \\
5 & Udalosť 2B & Sports         & Gymnastics      & 2.84 km & 5.0 (3)  & 2/74 \\
\hline
\end{tabular}
\caption{Vstupné údaje pre TOP 5 odporúčaní - Test 4}
\label{tab:data_test4}
\end{table}

V procese normalizácie (tabuľka \ref{tab:norm_test4}) dosahujú hodnoty $D$ (Vzdialenosť) u týchto udalostí maximálne úrovne, čo sa pri 60 \% váhe dramaticky prejavuje na celkovom výsledku.

\begin{table}[H]
\centering
\small
\begin{tabular}{|c|l|c|c|c|c|c|c|}
\hline
\textbf{\#} & \textbf{Event} & $C_{main}$ & $C_{sub}$ & $D$ & $R$ & $P$ & $B_{fav}$ \\
\hline
1 & Udalosť 2A & 0.375 & 0.000 & 0.879 & 0.94 & 0.714 & 0.0 \\
2 & Udalosť 5  & 0.375 & 0.000 & 0.863 & 0.88 & 0.147 & 0.0 \\
3 & Udalosť 1  & 0.050 & 0.000 & 0.905 & 0.94 & 0.061 & 0.0 \\
4 & Udalosť 8  & 0.500 & 0.000 & 0.797 & 0.68 & 0.294 & 0.0 \\
5 & Udalosť 2B & 0.500 & 0.000 & 0.779 & 1.00 & 0.027 & 0.0 \\
\hline
\end{tabular}
\caption{Normalizované hodnoty - Test 4}
\label{tab:norm_test4}
\end{table}

Detailný rozpis skóre v tabuľke \ref{tab:vysledky_test4} potvrdzuje, že faktor $w_3 D$ tvorí drvivú väčšinu (viac ako 45 \%) celkovej hodnoty skóre. Vďaka tomu môže byť udalosť s ratingom 3.4 (Udalosť 8) umiestnená vyššie ako kvalitnejšie, ale vzdialenejšie podujatia.

\begin{table}[H]
\centering
\small
\begin{tabular}{|c|l|c|c|c|c|c|c|c|}
\hline
\textbf{\#} & \textbf{Event} & \textbf{Score} & $w_1 C$ & $w_2 C_s$ & $w_3 D$ & $w_4 R$ & $w_5 P$ & $w_6 B$ \\
\hline
1 & Udalosť 2A & 68.9\% & 3.0\% & 0.0\% & 52.8\% & 7.5\% & 5.7\% & 0.0\% \\
2 & Udalosť 5  & 63.0\% & 3.0\% & 0.0\% & 51.7\% & 7.1\% & 1.2\% & 0.0\% \\
3 & Udalosť 1  & 62.8\% & 0.4\% & 0.0\% & 54.3\% & 7.6\% & 0.5\% & 0.0\% \\
4 & Udalosť 8  & 59.6\% & 4.0\% & 0.0\% & 47.8\% & 5.5\% & 2.4\% & 0.0\% \\
5 & Udalosť 2B & 58.9\% & 4.0\% & 0.0\% & 46.7\% & 7.9\% & 0.2\% & 0.0\% \\
\hline
\end{tabular}
\caption{Rozbor príspevkov k finálnemu skóre - Test 4}
\label{tab:vysledky_test4}
\end{table}

Nasledujúce grafy vizuálne ilustrujú totálnu dominanciu faktora vzdialenosti v tomto teste.

\begin{figure}[H]
\centering
\begin{minipage}{0.48\textwidth}
\centering
\begin{tikzpicture}
\begin{axis}[
    xbar, width=6cm, height=5cm,
    xlabel={Príspevok (\%)},
    symbolic y coords={Bonus, Popularita, Rating, Vzdialenosť, Podkat., Hl. kat.},
    ytick=data, nodes near coords, nodes near coords align={horizontal},
    title={Udalosť 2A - 68.9\%},
    xmin=0, xmax=65, bar width=0.25cm,
]
\addplot coordinates {(3.0,Hl. kat.) (0.0,Podkat.) (52.8,Vzdialenosť) (7.5,Rating) (5.7,Popularita) (0.0,Bonus)};
\end{axis}
\end{tikzpicture}
\end{minipage}
\hfill
\begin{minipage}{0.48\textwidth}
\centering
\begin{tikzpicture}
\begin{axis}[
    xbar, width=6cm, height=5cm,
    xlabel={Príspevok (\%)},
    symbolic y coords={Bonus, Popularita, Rating, Vzdialenosť, Podkat., Hl. kat.},
    ytick=data, nodes near coords, nodes near coords align={horizontal},
    title={Udalosť 5 - 63.0\%},
    xmin=0, xmax=65, bar width=0.25cm,
]
\addplot coordinates {(3.0,Hl. kat.) (0.0,Podkat.) (51.7,Vzdialenosť) (7.1,Rating) (1.2,Popularita) (0.0,Bonus)};
\end{axis}
\end{tikzpicture}
\end{minipage}
\caption{Vizuálny rozbor skóre pre Test 4 (Udalosť 2A a 5)}
\end{figure}

\begin{figure}[H]
\centering
\begin{minipage}{0.48\textwidth}
\centering
\begin{tikzpicture}
\begin{axis}[
    xbar, width=6cm, height=5cm,
    xlabel={Príspevok (\%)},
    symbolic y coords={Bonus, Popularita, Rating, Vzdialenosť, Podkat., Hl. kat.},
    ytick=data, nodes near coords, nodes near coords align={horizontal},
    title={Udalosť 1 - 62.8\%},
    xmin=0, xmax=65, bar width=0.25cm,
]
\addplot coordinates {(0.4,Hl. kat.) (0.0,Podkat.) (54.3,Vzdialenosť) (7.6,Rating) (0.5,Popularita) (0.0,Bonus)};
\end{axis}
\end{tikzpicture}
\end{minipage}
\hfill
\begin{minipage}{0.48\textwidth}
\centering
\begin{tikzpicture}
\begin{axis}[
    xbar, width=6cm, height=5cm,
    xlabel={Príspevok (\%)},
    symbolic y coords={Bonus, Popularita, Rating, Vzdialenosť, Podkat., Hl. kat.},
    ytick=data, nodes near coords, nodes near coords align={horizontal},
    title={Udalosť 8 - 59.6\%},
    xmin=0, xmax=65, bar width=0.25cm,
]
\addplot coordinates {(4.0,Hl. kat.) (0.0,Podkat.) (47.8,Vzdialenosť) (5.5,Rating) (2.4,Popularita) (0.0,Bonus)};
\end{axis}
\end{tikzpicture}
\end{minipage}
\caption{Vizuálny rozbor skóre pre Test 4 (Udalosť 1 a 8)}
\end{figure}

\begin{figure}[H]
\centering
\begin{tikzpicture}
\begin{axis}[
    xbar, width=10cm, height=5cm,
    xlabel={Príspevok k skóre (\%)},
    symbolic y coords={Bonus, Popularita, Rating, Vzdialenosť, Podkat., Hl. kat.},
    ytick=data, nodes near coords, nodes near coords align={horizontal},
    title={Udalosť 2B - 58.9\%},
    xmin=0, xmax=65, bar width=0.4cm,
]
\addplot coordinates {(4.0,Hl. kat.) (0.0,Podkat.) (46.7,Vzdialenosť) (7.9,Rating) (0.2,Popularita) (0.0,Bonus)};
\end{axis}
\end{tikzpicture}
\caption{Vizuálny rozbor skóre pre Test 4 (Udalosť 2B)}
\end{figure}



% ═══════════════════════════════════════════════════════════════════════════════
\subsection{Test 5: Prioritizácia používateľského hodnotenia (Rating 60 \%)}
% ═══════════════════════════════════════════════════════════════════════════════

V piatom scenári sme sa zamerali na situáciu, kedy je pre používateľa prvoradým kritériom kvalita podujatia vyjadrená priemerným hodnotením ostatných účastníkov. Váha faktora $w_4$ (Rating) bola zvýšená na 60 \%, zatiaľ čo ostatné faktory boli rovnomerne potlačené na úroveň 8 \%, ako uvádza tabuľka \ref{tab:vahy_test5}.

\begin{table}[H]
\centering
\begin{tabular}{|l|c|}
\hline
\textbf{Faktor} & \textbf{Váha ($w_i$)} \\
\hline
Hlavná kategória ($w_1$) & 8.00\% \\
Podkategória ($w_2$) & 8.00\% \\
Vzdialenosť ($w_3$) & 8.00\% \\
Rating ($w_4$) & 60.00\% \\
Popularita ($w_5$) & 8.00\% \\
Obľúbený bonus ($w_6$) & 8.00\% \\
\hline
\textbf{Súčet} & \textbf{100\%} \\
\hline
\end{tabular}
\caption{Konfigurácia váh pre Test 5}
\label{tab:vahy_test5}
\end{table}

V tabuľke \ref{tab:data_test5} sú uvedené základné parametre piatich najlepšie hodnotených podujatí. Je dôležité si všimnúť, že pri tejto konfigurácii už vzdialenosť (napr. 13 alebo 15 km) nehrá takú významnú rolu ako v predchádzajúcich testoch.

\begin{table}[H]
\centering
\small
\begin{tabular}{|c|l|l|l|c|c|c|}
\hline
\textbf{\#} & \textbf{Názov udalosti} & \textbf{Kategória} & \textbf{Podkategória} & \textbf{Vzdial.} & \textbf{Rating} & \textbf{Popularita} \\
\hline
1 & Udalosť 5  & Sports    & Strength \& Cond.   & 13.36 km & 4.9 (13) & 9/16 \\
2 & Udalosť 4A & Sports    & Team sports         & 15.80 km & 4.9 (1)  & 6/9 \\
3 & Udalosť 2A & Education & Jazykový kurz       & 1.37 km  & 4.7 (10) & 5/7 \\
4 & Udalosť 2B & Sports    & Gymnastics          & 2.84 km  & 5.0 (3)  & 2/74 \\
5 & Udalosť 3  & Sports    & Fitness \& Wellness & 3.73 km  & 4.9 (6)  & 9/96 \\
\hline
\end{tabular}
\caption{Vstupné údaje pre TOP 5 odporúčaní - Test 5}
\label{tab:data_test5}
\end{table}

\subsubsection{Normalizované hodnoty (bez váh)}

Pred aplikáciou 60 \% váhy na rating systém prepočíta surové dáta na bezrozmerné koeficienty (tabuľka \ref{tab:norm_test5}). Všimnite si, že pri udalosti 4 dosahuje faktor $R$ (Rating) maximálnu hodnotu 1.000, čo zodpovedá čistému 5.0 hodnoteniu.

\begin{table}[H]
\centering
\small
\begin{tabular}{|c|l|c|c|c|c|c|c|}
\hline
\textbf{\#} & \textbf{Event} & $C_{main}$ & $C_{sub}$ & $D$ & $R$ & $P$ & $B_{fav}$ \\
\hline
1 & Udalosť 5  & 0.500 & 0.125 & 0.428 & 0.98 & 0.563 & 1.0 \\
2 & Udalosť 4A & 0.500 & 0.125 & 0.387 & 0.98 & 0.667 & 0.0 \\
3 & Udalosť 2A & 0.375 & 0.000 & 0.879 & 0.94 & 0.714 & 0.0 \\
4 & Udalosť 2B & 0.500 & 0.000 & 0.779 & 1.00 & 0.027 & 0.0 \\
5 & Udalosť 3  & 0.500 & 0.125 & 0.728 & 0.98 & 0.094 & 0.0 \\
\hline
\end{tabular}
\caption{Normalizované vstupné hodnoty - Test 5}
\label{tab:norm_test5}
\end{table}

Finálne vypočítané skóre (tabuľka \ref{tab:vysledky_test5}) jasne dominuje hodnota Ratingu. Aj pri udalosti 4, ktorá má len jedno hodnotenie (100 % úspešnosť), dokáže algoritmus správne identifikovať vysoký potenciál spokojnosti.

\begin{table}[H]
\centering
\small
\begin{tabular}{|c|l|c|c|c|c|c|c|c|}
\hline
\textbf{\#} & \textbf{Event} & \textbf{Score} & $w_1 C$ & $w_2 C_s$ & $w_3 D$ & $w_4 R$ & $w_5 P$ & $w_6 B$ \\
\hline
1 & Udalosť 5  & 77.1\% & 4.0\% & 1.0\% & 3.4\% & 58.8\% & 4.5\% & 8.0\% \\
2 & Udalosť 4A & 69.1\% & 4.0\% & 1.0\% & 3.1\% & 58.8\% & 5.3\% & 0.0\% \\
3 & Udalosť 2A & 67.2\% & 3.0\% & 0.0\% & 7.0\% & 56.4\% & 5.7\% & 0.0\% \\
4 & Udalosť 2B & 66.8\% & 4.0\% & 0.0\% & 6.2\% & 60.0\% & 0.2\% & 0.0\% \\
5 & Udalosť 3  & 66.4\% & 4.0\% & 1.0\% & 5.8\% & 58.8\% & 0.8\% & 0.0\% \\
\hline
\end{tabular}
\caption{Rozbor príspevkov k finálnemu skóre - Test 5}
\label{tab:vysledky_test5}
\end{table}

Vizuálna analýza na obrázkoch nižšie potvrdzuje, že faktor Ratingu (tretí stĺpec zospodu) tvorí pri každom podujatí dominantnú časť celkového výsledku.

\begin{figure}[H]
\centering
\begin{minipage}{0.48\textwidth}
\centering
\begin{tikzpicture}
\begin{axis}[
    xbar, width=6cm, height=5cm,
    xlabel={Príspevok (\%)},
    symbolic y coords={Bonus, Popularita, Rating, Vzdialenosť, Podkat., Hl. kat.},
    ytick=data, nodes near coords, nodes near coords align={horizontal},
    title={Udalosť 5 - 77.1\%},
    xmin=0, xmax=65, bar width=0.25cm,
]
\addplot coordinates {(4.0,Hl. kat.) (1.0,Podkat.) (3.4,Vzdialenosť) (58.8,Rating) (4.5,Popularita) (8.0,Bonus)};
\end{axis}
\end{tikzpicture}
\end{minipage}
\hfill
\begin{minipage}{0.48\textwidth}
\centering
\begin{tikzpicture}
\begin{axis}[
    xbar, width=6cm, height=5cm,
    xlabel={Príspevok (\%)},
    symbolic y coords={Bonus, Popularita, Rating, Vzdialenosť, Podkat., Hl. kat.},
    ytick=data, nodes near coords, nodes near coords align={horizontal},
    title={Udalosť 4A - 69.1\%},
    xmin=0, xmax=65, bar width=0.25cm,
]
\addplot coordinates {(4.0,Hl. kat.) (1.0,Podkat.) (3.1,Vzdialenosť) (58.8,Rating) (5.3,Popularita) (0.0,Bonus)};
\end{axis}
\end{tikzpicture}
\end{minipage}
\caption{Vizuálny rozbor skóre pre Test 5 (Udalosť 5 a 4A)}
\end{figure}

\begin{figure}[H]
\centering
\begin{minipage}{0.48\textwidth}
\centering
\begin{tikzpicture}
\begin{axis}[
    xbar, width=6cm, height=5cm,
    xlabel={Príspevok (\%)},
    symbolic y coords={Bonus, Popularita, Rating, Vzdialenosť, Podkat., Hl. kat.},
    ytick=data, nodes near coords, nodes near coords align={horizontal},
    title={Udalosť 2A - 67.2\%},
    xmin=0, xmax=65, bar width=0.25cm,
]
\addplot coordinates {(3.0,Hl. kat.) (0.0,Podkat.) (7.0,Vzdialenosť) (56.4,Rating) (5.7,Popularita) (0.0,Bonus)};
\end{axis}
\end{tikzpicture}
\end{minipage}
\hfill
\begin{minipage}{0.48\textwidth}
\centering
\begin{tikzpicture}
\begin{axis}[
    xbar, width=6cm, height=5cm,
    xlabel={Príspevok (\%)},
    symbolic y coords={Bonus, Popularita, Rating, Vzdialenosť, Podkat., Hl. kat.},
    ytick=data, nodes near coords, nodes near coords align={horizontal},
    title={Udalosť 2B - 66.8\%},
    xmin=0, xmax=65, bar width=0.25cm,
]
\addplot coordinates {(4.0,Hl. kat.) (0.0,Podkat.) (6.2,Vzdialenosť) (60.0,Rating) (0.2,Popularita) (0.0,Bonus)};
\end{axis}
\end{tikzpicture}
\end{minipage}
\caption{Vizuálny rozbor skóre pre Test 5 (Udalosť 2A and 2B)}
\end{figure}

\begin{figure}[H]
\centering
\begin{tikzpicture}
\begin{axis}[
    xbar, width=10cm, height=5cm,
    xlabel={Príspevok k skóre (\%)},
    symbolic y coords={Bonus, Popularita, Rating, Vzdialenosť, Podkat., Hl. kat.},
    ytick=data, nodes near coords, nodes near coords align={horizontal},
    title={Udalosť 3 - 66.4\%},
    xmin=0, xmax=65, bar width=0.4cm,
]
\addplot coordinates {(4.0,Hl. kat.) (1.0,Podkat.) (5.8,Vzdialenosť) (58.8,Rating) (0.8,Popularita) (0.0,Bonus)};
\end{axis}
\end{tikzpicture}
\caption{Vizuálny rozbor skóre pre Test 5 (Udalosť 3)}
\end{figure}


% ═══════════════════════════════════════════════════════════════════════════════
% ═══════════════════════════════════════════════════════════════════════════════
\subsection{Test 6: Prioritizácia popularity a obsadenosti (Popularita 60 \%)}
% ═══════════════════════════════════════════════════════════════════════════════

V šiestom scenári simulujeme správanie systému pre používateľa, ktorý vyhľadáva spoločensky exponované a populárne podujatia (tzv. "trending events"). Váha faktora $w_5$ (Popularita), ktorý je v systéme definovaný ako pomer prihlásených účastníkov k celkovej kapacite, bola nastavená na 60 \%. Ostatné váhy sú rovnomerne rozdelené po 8 \%, ako vidieť v tabuľke \ref{tab:vahy_test6}.

\begin{table}[H]
\centering
\begin{tabular}{|l|c|}
\hline
\textbf{Faktor} & \textbf{Váha ($w_i$)} \\
\hline
Hlavná kategória ($w_1$) & 8.00\% \\
Podkategória ($w_2$) & 8.00\% \\
Vzdialenosť ($w_3$) & 8.00\% \\
Rating ($w_4$) & 8.00\% \\
Popularita ($w_5$) & 60.00\% \\
Obľúbený bonus ($w_6$) & 8.00\% \\
\hline
\textbf{Súčet} & \textbf{100\%} \\
\hline
\end{tabular}
\caption{Konfigurácia váh pre Test 6}
\label{tab:vahy_test6}
\end{table}

Vstupné údaje (tabuľka \ref{tab:data_test6}) ukazujú vysokú mieru naplnenia kapacity pri všetkých piatich udalostiach (napr. 5/7 pri Udalosti 2 alebo 12/15 pri Udalosti 11). To spôsobuje, že sa do výberu dostávajú podujatia z rôznych kategórií bez ohľadu na ich vzájomnú tematickú súvislosť.

\begin{table}[H]
\centering
\small
\begin{tabular}{|c|l|l|l|c|c|c|}
\hline
\textbf{\#} & \textbf{Názov udalosti} & \textbf{Kategória} & \textbf{Podkategória} & \textbf{Vzdial.} & \textbf{Rating} & \textbf{Popularita} \\
\hline
1 & Udalosť 2A & Education          & Jazykový kurz  & 1.37 km  & 4.7 (10) & 5/7 \\
2 & Udalosť 11 & Community \& Soc.  & Speed dating   & 9.18 km  & 4.2 (16) & 12/15 \\
3 & Udalosť 7  & Food \& Drinks     & Food festival  & 20.71 km & 3.9 (7)  & 9/11 \\
4 & Udalosť 5  & Sports             & Strength \& C. & 13.36 km & 4.9 (13) & 9/16 \\
5 & Udalosť 4B & Sports             & Team sports    & 15.80 km & 4.9 (1)  & 6/9 \\
\hline
\end{tabular}
\caption{Vstupné údaje pre TOP 5 odporúčaní - Test 6}
\label{tab:data_test6}
\end{table}

Normalizované hodnoty (tabuľka \ref{tab:norm_test6}) potvrdzujú, že faktor $P$ (Popularita) dosahuje v tomto teste najvyššie relatívne hodnoty v porovnaní s ostatnými atribútmi.

\begin{table}[H]
\centering
\small
\begin{tabular}{|c|l|c|c|c|c|c|c|}
\hline
\textbf{\#} & \textbf{Event} & $C_{main}$ & $C_{sub}$ & $D$ & $R$ & $P$ & $B_{fav}$ \\
\hline
1 & Udalosť 2A & 0.375 & 0.000 & 0.879 & 0.94 & 0.714 & 0.0 \\
2 & Udalosť 11 & 0.050 & 0.000 & 0.521 & 0.84 & 0.800 & 0.0 \\
3 & Udalosť 7  & 0.050 & 0.000 & 0.326 & 0.78 & 0.818 & 0.0 \\
4 & Udalosť 5  & 0.500 & 0.125 & 0.428 & 0.98 & 0.563 & 1.0 \\
5 & Udalosť 4B & 0.500 & 0.125 & 0.387 & 0.98 & 0.667 & 0.0 \\
\hline
\end{tabular}
\caption{Normalizované hodnoty - Test 6}
\label{tab:norm_test6}
\end{table}



Výsledné skóre (tabuľka \ref{tab:vysledky_test6}) ukazuje, že Popularita ($w_5 P$) tvorí u najúspešnejšej udalosti takmer 43 \% z celkových 60,4 \%. Tento scenár demonštruje, že systém dokáže efektívne "vytiahnuť" akcie, o ktoré je v danej chvíli najväčší záujem.

\begin{table}[H]
\centering
\small
\begin{tabular}{|c|l|c|c|c|c|c|c|c|}
\hline
\textbf{\#} & \textbf{Event} & \textbf{Score} & $w_1 C$ & $w_2 C_s$ & $w_3 D$ & $w_4 R$ & $w_5 P$ & $w_6 B$ \\
\hline
1 & Udalosť 2A & 60.4\% & 3.0\% & 0.0\% & 7.0\% & 7.5\% & 42.9\% & 0.0\% \\
2 & Udalosť 11 & 59.3\% & 0.4\% & 0.0\% & 4.2\% & 6.8\% & 48.0\% & 0.0\% \\
3 & Udalosť 7  & 58.3\% & 0.4\% & 0.0\% & 2.6\% & 6.2\% & 49.1\% & 0.0\% \\
4 & Udalosť 5  & 58.1\% & 4.0\% & 1.0\% & 3.4\% & 7.9\% & 33.8\% & 8.0\% \\
5 & Udalosť 4B & 56.0\% & 4.0\% & 1.0\% & 3.1\% & 7.9\% & 40.0\% & 0.0\% \\
\hline
\end{tabular}
\caption{Rozbor príspevkov k finálnemu skóre - Test 6}
\label{tab:vysledky_test6}
\end{table}

Vizualizácia nižšie potvrdzuje dominanciu faktora popularity (stĺpec Popularita) u všetkých zobrazených udalostí.

\begin{figure}[H]
\centering
\begin{minipage}{0.48\textwidth}
\centering
\begin{tikzpicture}
\begin{axis}[
    xbar, width=6cm, height=5cm,
    xlabel={Príspevok (\%)},
    symbolic y coords={Bonus, Popularita, Rating, Vzdialenosť, Podkat., Hl. kat.},
    ytick=data, nodes near coords, nodes near coords align={horizontal},
    title={Udalosť 2A - 60.4\%},
    xmin=0, xmax=55, bar width=0.25cm,
]
\addplot coordinates {(3.0,Hl. kat.) (0.0,Podkat.) (7.0,Vzdialenosť) (7.5,Rating) (42.9,Popularita) (0.0,Bonus)};
\end{axis}
\end{tikzpicture}
\end{minipage}
\hfill
\begin{minipage}{0.48\textwidth}
\centering
\begin{tikzpicture}
\begin{axis}[
    xbar, width=6cm, height=5cm,
    xlabel={Príspevok (\%)},
    symbolic y coords={Bonus, Popularita, Rating, Vzdialenosť, Podkat., Hl. kat.},
    ytick=data, nodes near coords, nodes near coords align={horizontal},
    title={Udalosť 11 - 59.3\%},
    xmin=0, xmax=55, bar width=0.25cm,
]
\addplot coordinates {(0.4,Hl. kat.) (0.0,Podkat.) (4.2,Vzdialenosť) (6.8,Rating) (48.0,Popularita) (0.0,Bonus)};
\end{axis}
\end{tikzpicture}
\end{minipage}
\caption{Vizuálny rozbor skóre pre Test 6 (Udalosť 2A a 11)}
\end{figure}

\begin{figure}[H]
\centering
\begin{minipage}{0.48\textwidth}
\centering
\begin{tikzpicture}
\begin{axis}[
    xbar, width=6cm, height=5cm,
    xlabel={Príspevok (\%)},
    symbolic y coords={Bonus, Popularita, Rating, Vzdialenosť, Podkat., Hl. kat.},
    ytick=data, nodes near coords, nodes near coords align={horizontal},
    title={Udalosť 7 - 58.3\%},
    xmin=0, xmax=55, bar width=0.25cm,
]
\addplot coordinates {(0.4,Hl. kat.) (0.0,Podkat.) (2.6,Vzdialenosť) (6.2,Rating) (49.1,Popularita) (0.0,Bonus)};
\end{axis}
\end{tikzpicture}
\end{minipage}
\hfill
\begin{minipage}{0.48\textwidth}
\centering
\begin{tikzpicture}
\begin{axis}[
    xbar, width=6cm, height=5cm,
    xlabel={Príspevok (\%)},
    symbolic y coords={Bonus, Popularita, Rating, Vzdialenosť, Podkat., Hl. kat.},
    ytick=data, nodes near coords, nodes near coords align={horizontal},
    title={Udalosť 5 - 58.1\%},
    xmin=0, xmax=55, bar width=0.25cm,
]
\addplot coordinates {(4.0,Hl. kat.) (1.0,Podkat.) (3.4,Vzdialenosť) (7.9,Rating) (33.8,Popularita) (8.0,Bonus)};
\end{axis}
\end{tikzpicture}
\end{minipage}
\caption{Vizuálny rozbor skóre pre Test 6 (Udalosť 7 a 5)}
\end{figure}

\begin{figure}[H]
\centering
\begin{tikzpicture}
\begin{axis}[
    xbar, width=10cm, height=5cm,
    xlabel={Príspevok k skóre (\%)},
    symbolic y coords={Bonus, Popularita, Rating, Vzdialenosť, Podkat., Hl. kat.},
    ytick=data, nodes near coords, nodes near coords align={horizontal},
    title={Udalosť 4B - 56.0\%},
    xmin=0, xmax=55, bar width=0.4cm,
]
\addplot coordinates {(4.0,Hl. kat.) (1.0,Podkat.) (3.1,Vzdialenosť) (7.9,Rating) (40.0,Popularita) (0.0,Bonus)};
\end{axis}
\end{tikzpicture}
\caption{Vizuálny rozbor skóre pre Test 6 (Udalosť 4B)}
\end{figure}

% ═══════════════════════════════════════════════════════════════════════════════
\subsection{Test 7: Prioritizácia obľúbených položiek (Obľúbený bonus 60 \%)}
% ═══════════════════════════════════════════════════════════════════════════════

Záverečný testovací scenár sa zameriava na personalizáciu prostredníctvom faktora $w_6$ (Obľúbený bonus). Tento atribút nadobúda hodnotu 1, ak udalosť organizuje subjekt, ktorého má používateľ v zozname obľúbených, inak je nulový. Nastavením váhy na 60 \% (tabuľka \ref{tab:vahy_test7}) simulujeme prostredie, kde systém prioritne doručuje obsah od overených organizátorov.

\begin{table}[H]
\centering
\begin{tabular}{|l|c|}
\hline
\textbf{Faktor} & \textbf{Váha ($w_i$)} \\
\hline
Hlavná kategória ($w_1$) & 8.00\% \\
Podkategória ($w_2$) & 8.00\% \\
Vzdialenosť ($w_3$) & 8.00\% \\
Rating ($w_4$) & 8.00\% \\
Popularita ($w_5$) & 8.00\% \\
Obľúbený bonus ($w_6$) & 60.00\% \\
\hline
\textbf{Súčet} & \textbf{100\%} \\
\hline
\end{tabular}
\caption{Konfigurácia váh pre Test 7}
\label{tab:vahy_test7}
\end{table}

Vstupná množina dát (tabuľka \ref{tab:data_test7}) ukazuje, že prvé tri udalosti majú priradený bonus za obľúbeného organizátora. Hoci majú udalosti 4 a 5 v niektorých metrikách (napr. vzdialenosť alebo popularita) lepšie výsledky, absencia bonusu ich v konečnom poradí výrazne odsunula.

\begin{table}[H]
\centering
\small
\begin{tabular}{|c|l|l|l|c|c|c|}
\hline
\textbf{\#} & \textbf{Názov udalosti} & \textbf{Kategória} & \textbf{Podkategória} & \textbf{Vzdial.} & \textbf{Rating} & \textbf{Popularita} \\
\hline
1 & Udalosť 5  & Sports    & Strength \& Cond. & 13.36 km & 4.9 (13) & 9/16 \\
2 & Udalosť 15 & Sports    & Strength \& Cond. & 11.73 km & 4.3 (17) & 4/71 \\
3 & Udalosť 7  & Sports    & Strength \& Cond. & 15.63 km & 4.2 (14) & 6/69 \\
4 & Udalosť 2  & Education & Jazykový kurz     & 1.37 km  & 4.7 (10) & 5/7 \\
5 & Udalosť 4  & Sports    & Team sports       & 15.80 km & 4.9 (1)  & 6/9 \\
\hline
\end{tabular}
\caption{Vstupné údaje pre TOP 5 odporúčaní - Test 7}
\label{tab:data_test7}
\end{table}



Normalizované hodnoty v tabuľke \ref{tab:norm_test7} jasne identifikujú prítomnosť binárneho bonusu $B_{fav}$ pre prvé tri položky.

\begin{table}[H]
\centering
\small
\begin{tabular}{|c|l|c|c|c|c|c|c|}
\hline
\textbf{\#} & \textbf{Event} & $C_{main}$ & $C_{sub}$ & $D$ & $R$ & $P$ & $B_{fav}$ \\
\hline
1 & Udalosť 5  & 0.500 & 0.125 & 0.428 & 0.98 & 0.563 & 1.0 \\
2 & Udalosť 15 & 0.500 & 0.125 & 0.460 & 0.86 & 0.056 & 1.0 \\
3 & Udalosť 7  & 0.500 & 0.125 & 0.390 & 0.84 & 0.087 & 1.0 \\
4 & Udalosť 2  & 0.375 & 0.000 & 0.879 & 0.94 & 0.714 & 0.0 \\
5 & Udalosť 4  & 0.500 & 0.125 & 0.387 & 0.98 & 0.667 & 0.0 \\
\hline
\end{tabular}
\caption{Normalizované hodnoty - Test 7}
\label{tab:norm_test7}
\end{table}

Výsledné skóre (tabuľka \ref{tab:vysledky_test7}) odhaľuje zásadný rozdiel v bodovom hodnotení. Zatiaľ čo udalosti s bonusom prekračujú hranicu 75 \%, najlepšia udalosť bez bonusu (Udalosť 2) dosahuje len 23,2 \%. Tento test potvrdzuje, že váha 60 \% pre binárny faktor pôsobí v algoritme ako dominantný prioritizačný filter.

\begin{table}[H]
\centering
\small
\begin{tabular}{|c|l|c|c|c|c|c|c|c|}
\hline
\textbf{\#} & \textbf{Event} & \textbf{Score} & $w_1 C$ & $w_2 C_s$ & $w_3 D$ & $w_4 R$ & $w_5 P$ & $w_6 B$ \\
\hline
1 & Udalosť 5  & 80.8\% & 4.0\% & 1.0\% & 3.4\% & 7.9\% & 4.5\% & 60.0\% \\
2 & Udalosť 15 & 76.1\% & 4.0\% & 1.0\% & 3.7\% & 6.9\% & 0.5\% & 60.0\% \\
3 & Udalosť 7  & 75.5\% & 4.0\% & 1.0\% & 3.1\% & 6.7\% & 0.7\% & 60.0\% \\
4 & Udalosť 2  & 23.2\% & 3.0\% & 0.0\% & 7.0\% & 7.5\% & 5.7\% & 0.0\% \\
5 & Udalosť 4  & 21.3\% & 4.0\% & 1.0\% & 3.1\% & 7.9\% & 5.3\% & 0.0\% \\
\hline
\end{tabular}
\caption{Rozbor príspevkov k finálnemu skóre - Test 7}
\label{tab:vysledky_test7}
\end{table}

Nasledujúce grafy vizualizujú tento priepastný rozdiel. Pri prvých troch udalostiach tvorí stĺpec "Bonus" drvivú väčšinu celkového skóre.

\begin{figure}[H]
\centering
\begin{minipage}{0.48\textwidth}
\centering
\begin{tikzpicture}
\begin{axis}[
    xbar, width=6cm, height=5cm,
    xlabel={Príspevok (\%)},
    symbolic y coords={Bonus, Popularita, Rating, Vzdialenosť, Podkat., Hl. kat.},
    ytick=data, nodes near coords, nodes near coords align={horizontal},
    title={Udalosť 5 - 80.8\%},
    xmin=0, xmax=65, bar width=0.25cm,
]
\addplot coordinates {(4.0,Hl. kat.) (1.0,Podkat.) (3.4,Vzdialenosť) (7.9,Rating) (4.5,Popularita) (60.0,Bonus)};
\end{axis}
\end{tikzpicture}
\end{minipage}
\hfill
\begin{minipage}{0.48\textwidth}
\centering
\begin{tikzpicture}
\begin{axis}[
    xbar, width=6cm, height=5cm,
    xlabel={Príspevok (\%)},
    symbolic y coords={Bonus, Popularita, Rating, Vzdialenosť, Podkat., Hl. kat.},
    ytick=data, nodes near coords, nodes near coords align={horizontal},
    title={Udalosť 15 - 76.1\%},
    xmin=0, xmax=65, bar width=0.25cm,
]
\addplot coordinates {(4.0,Hl. kat.) (1.0,Podkat.) (3.7,Vzdialenosť) (6.9,Rating) (0.5,Popularita) (60.0,Bonus)};
\end{axis}
\end{tikzpicture}
\end{minipage}
\caption{Vizuálny rozbor skóre pre Test 7 (Udalosť 5 a 15)}
\end{figure}

\begin{figure}[H]
\centering
\begin{minipage}{0.48\textwidth}
\centering
\begin{tikzpicture}
\begin{axis}[
    xbar, width=6cm, height=5cm,
    xlabel={Príspevok (\%)},
    symbolic y coords={Bonus, Popularita, Rating, Vzdialenosť, Podkat., Hl. kat.},
    ytick=data, nodes near coords, nodes near coords align={horizontal},
    title={Udalosť 7 - 75.5\%},
    xmin=0, xmax=65, bar width=0.25cm,
]
\addplot coordinates {(4.0,Hl. kat.) (1.0,Podkat.) (3.1,Vzdialenosť) (6.7,Rating) (0.7,Popularita) (60.0,Bonus)};
\end{axis}
\end{tikzpicture}
\end{minipage}
\hfill
\begin{minipage}{0.48\textwidth}
\centering
\begin{tikzpicture}
\begin{axis}[
    xbar, width=6cm, height=5cm,
    xlabel={Príspevok (\%)},
    symbolic y coords={Bonus, Popularita, Rating, Vzdialenosť, Podkat., Hl. kat.},
    ytick=data, nodes near coords, nodes near coords align={horizontal},
    title={Udalosť 2 - 23.2\%},
    xmin=0, xmax=65, bar width=0.25cm,
]
\addplot coordinates {(3.0,Hl. kat.) (0.0,Podkat.) (7.0,Vzdialenosť) (7.5,Rating) (5.7,Popularita) (0.0,Bonus)};
\end{axis}
\end{tikzpicture}
\end{minipage}
\caption{Vizuálny rozbor skóre pre Test 7 (Udalosť 7 a 2)}
\end{figure}

\begin{figure}[H]
\centering
\begin{tikzpicture}
\begin{axis}[
    xbar, width=10cm, height=5cm,
    xlabel={Príspevok k skóre (\%)},
    symbolic y coords={Bonus, Popularita, Rating, Vzdialenosť, Podkat., Hl. kat.},
    ytick=data, nodes near coords, nodes near coords align={horizontal},
    title={Udalosť 4 - 21.3\%},
    xmin=0, xmax=65, bar width=0.4cm,
]
\addplot coordinates {(4.0,Hl. kat.) (1.0,Podkat.) (3.1,Vzdialenosť) (7.9,Rating) (5.3,Popularita) (0.0,Bonus)};
\end{axis}
\end{tikzpicture}
\caption{Vizuálny rozbor skóre pre Test 7 (Udalosť 4)}
\end{figure}
\end{document}